\Needspace{18\baselineskip}
The preceding lemmas give a canonical ordinary type for a countable divisorial current of finite mass. Its envelope characterization uses the ordinary and ideal envelopes of \cref{def:singularity-envelope,def:hermitian-i-envelope}.
\begin{corollary}\label{cor:divisorial-current-ordinary-type}
Let $Z$ be a connected compact complex manifold, let $D_j$ be distinct prime divisors with integration currents $[D_j]$, and let $c_j>0$. Suppose $S=\sum_jc_j[D_j]$ is a closed positive current of finite mass. It determines a canonical ordinary quasi-plurisubharmonic type with representatives $v$ satisfying
\[
 \ddc v=S-\beta
\]
for a smooth closed real form $\beta$. Any two such representatives differ smoothly.
For every Hermitian background $\Omega$ admitting the chosen representative,
\[
 P_\Omega[v]=P_\Omega[v]_{\mathcal I}\sim v.
\]
\end{corollary}
\begin{proof}
\Cref{lem:closed-current-global-potential} supplies the decomposition and shows that its potentials differ smoothly. Thus the ordinary type depends only on $S$.

By \cref{thm:hermitian-i-projection}, the ideal envelope preserves all multiplier ideals, so its generic Lelong number on each $D_j$ is still $c_j$. \Cref{lem:divisorial-potential-domination} gives $P_\Omega[v]_{\mathcal I}\preceq v$. Also $P_\Omega[v]_{\mathcal I}\geq v-\sup v$, so these functions are ordinarily equivalent. Consequently $P_\Omega[v]_{\mathcal I}$ is admissible in the ordinary envelope of $v$, giving $P_\Omega[v]_{\mathcal I}\leq P_\Omega[v]$. The reverse inequality is \cref{prop:p-below-i}. This proves the envelope identity and ordinary equivalence.
\end{proof}
